Praktično vse naloge v mehaniki vsebujejo obravnavo gibanja nekega telesa.
Da to naredimo, analiziramo sile na telo, in zapišemo drugi Newtonov zakon
\[
  m \ddot{\vec{x}} = \sum \vec{F}.
\]
Ker je sila $\vec{F}$ lahko odvisna od časa $t$, položaja telesa $\vec{x}$ ter
od njegove hitrosti $\dot{\vec{x}}$, dobimo diferencialno enačbo drugega reda.
Te rešujemo z običajnimi triki za reševanje diferencialnih enačb.

Če za silo $\vec{F}$ obstaja funkcija $U: \R^3 \to \R$, da je $\vec{F} = - \grad
U$, pravimo, da je $\vec{F}$ \pojem{potencialna} ali \pojem{konzervativna} sila.
Delo take sile je odvisno le od začetnega in končnega položaja točke; velja
$W + U = E_0$, kjer je $W$ kinetična energija in $E_0$ neka konstanta.
Pri obravnavi premočrtnega gibanja pod vplivom potencialne sile naredimo dve
stvari.
Najprej narišemo graf danega potenciala $U(x)$, s katerim naredimo kvalitativno
analizo gibanja.
Če k grafu narišemo še horizontalno črto za $E_0$, lahko povemo, kaj se bo
zgodilo z delcem na tem energijskem nivoju;
\begin{itemize}
\item Če je $E_0$ popolnoma nad grafom $U(x)$, bo gibanje neomejeno.
  Tukaj lahko dodatno povemo, če dosežemo točko v neskončnosti v končnem ali
  neskončnem času, kar je odvisno od grafa $U(x)$; razlika med $E_0$ in $U(x)$
  pove kinetično energijo, ki jo bo delec imel v $x$.
\item Če $E_0$ seka graf $U(x)$ v dveh točkah, bo gibanje omejeno (in
  periodično).
  V tem primeru lahko izračunamo periodo gibanja, kot piše spodaj.
\item Če se $E_0$ dotika grafa v lokalnem maksimumu, bo gibanje v to smer
  omejeno, ker v točko dotika ne moramo priti v končnem času.
\item Če se $E_0$ dotika grafa v lokalnem minimumu, smo v stabilnem ravnovesju,
  in se ne bomo nikdar premaknili iz te točke.
\end{itemize}
Pri reševanju naloge zapišemo vse možnosti za gibanje v odvisnosti od $E_0$.

Pogosto moramo izračunati tudi periodo nihanja v primeru druge točke, kar lahko
naredimo direktno, ali pa z eno od dveh aproksimacij.
Recimo, da se gibljemo med točkama $a < b$.
Tedaj lahko periodo izračunamo kot
\[
  T = \sqrt{2m} \int_a^b \frac{dx}{\sqrt{E_0 - U(x)}}.
\]
To ni le integral, pač pa celo posplošen integral (ker je $U(x) = E_0$ v
krajiščih).
Sicer je vedno končen, a ga je pogosto težko (nemogoče) izračunati.
Pri računu si lahko pomagamo z integralom
\[
  \int_a^b \frac{dx}{\sqrt{(b-x)(x-a)}} = \pi,
\]
kjer pazimo, da sta vrednosti v ulomku enaki mejam integrala.

V posebnem primeru harmoničnega oscilatorja velja $U = \pol k x^2$ za neko
konstanto $k$.
Tedaj so vsa gibanja periodična s periodo $T = 2 \pi \sqrt{\nicefrac{m}{k}}$.
Če je perioda gibanja neodvisna od $E_0$, kakor je tu, pravimo, da je gibanje
\pojem{izhronično}.
Izkaže se, da je edini potencial, ki je hkrati izohroničen in simetričen glede
na svoj minimum, harmonični.
Naj bo potencial $U$ sedaj neharmonični.
Če v lokalnem minimumu $x_0$ velja $E_0 - U(x) \ll 1$, se lahko poslužimo
harmonične aproksimacije.
Pri njej se predstavljamo, da je potencial kvadratna funkcija v tej točki, in
izračunamo približek
\[
  T \doteq 2 \pi \sqrt{\frac{m}{U''(x_0)}}.
\]

Druga možna aproksimacija je \pojem{libracijska}, kjer prvo zapišemo
\[
  E_0 - U(x) = (b-x)(x-a) \chi(x),
\]
in ocenimo integral po trapezni formuli, tako da dobimo približek
\[
  T \doteq \pi \sqrt{\frac{m}{2}} \left( \inv{\sqrt{\chi(a)}} +
	\inv{\sqrt{\chi(b)}} \right),
\]
kar je natančna ocena v primeru, da je funkcija ${(\chi(z))}^{-1/2}$ afina.